\documentclass[10pt,a4paper]{amsart}
\usepackage{amsmath,amssymb,amsthm,graphicx}
\usepackage{hyperref}
\numberwithin{equation}{section}

\newtheorem{Theorem}{Theorem}
\newtheorem*{Theorem*}{Theorem}
\newtheorem{Corollary}[Theorem]{Corollary}
\newtheorem{Lemma}[Theorem]{Lemma}
\newtheorem{Proposition}[Theorem]{Proposition}

\theoremstyle{definition}
\newtheorem{Definition}[Theorem]{Definition}
\newtheorem{Note}[Theorem]{Note}
\newtheorem{Example}[Theorem]{Example}
\newtheorem{Remark}[Theorem]{Remark}

\usepackage{tikz-cd}

\title{Complements of Reeb-invariant immaterial sets surject onto symplectic homology}
\author{Igor Uljarevic}
\date{doi:10.3842/SIGMA.2026.051; CC BY 4.0}
\begin{document}
\maketitle
\noindent\textit{Allowed standard prerequisites (editorial).}\begin{itemize}
\item Standard contact/Floer continuation and naturality. Original complete bibliography and source-local definitions retained.
\end{itemize}\medskip

A subset $A$ of a cooriented contact manifold $M$ is called \emph{immaterial}\footnote{The notion was introduced in \cite[Section~7]{uljarevic2023selective}.} if there exists a contractible loop of contactomorphisms $\varphi_t\colon M\to M$ whose vector field $X$ points in the positive direction (of the contact distribution) on $A$.
The property of being immaterial is clearly related to orderability, a well studied phenomenon in contact geometry originally introduced in \cite{eliashberg2000partially}. Namely, a contact manifold~$M$ is non-orderable if, and only if, all of its subsets are immaterial. The existence of non-orderable contact manifolds shows that immaterial subsets are not necessarily small in the sense of contact squeezing or contact displaceability. Indeed, the whole contact manifold~$M$ (within itself) cannot be squeezed into a smaller subset or displaced from itself, although the whole $M$ is immaterial if non-orderable.
The only immaterial subset in dimension 1 is the empty set.\footnote{{Here, we prove the claim for $\mathbb{S}^1$. It is enough to show that a point $p\in\mathbb{S}^1$ is not immaterial. Let $\varphi_t\colon \mathbb{S}^1\to \mathbb{S}^1$ be a loop of contactomorphisms that is positive at $p$. Then, the degree of the map $\gamma\colon\mathbb{S}^1\to \mathbb{S}^1$, $t\mapsto \varphi_t(p)$ is positive. (This can be seen by counting elements of $\gamma^{-1}(\{p\})$ with signs.) Therefore, $\{\varphi_t\}_t$ cannot be contractible.}} This is in contrast with dimensions greater than 3 where every compact subset of a Darboux chart is immaterial \cite[Lemma~7.3]{uljarevic2023selective}.

{The results of this note can be seen as contributions to questions regarding the size in contact geometry. These questions are quite delicate (as evident from work by Eliashberg--Kim--Polterovich \cite{eliashberg2006geometry}) and have been an active research area over the last two decades. Our first result (Theorem~\ref{thm:complement-big}) is formulated in terms of selective symplectic homology \cite{uljarevic2023selective}, while the second (Theorem~\ref{thm:tau}) is expressed using recently introduced contact quasi-measures \cite{djordjevic2025quantitative}. The contact quasi-measures from \cite{djordjevic2025quantitative} are strongly inspired by Entov--Polterovich's construction of (symplectic) quasi-measures \cite{EP06} using spectral invariants from \cite{oh2004construction}. See \cite{Ent14} or \cite{polterovich2014function} for a more detailed historical account.}

{A central role in this paper is played by \emph{contact Hamiltonian Floer homology}, introduced in~\cite{merry2019maximum} by proving a maximum principle for slopes that are allowed to be general contact Hamiltonians (instead of constants in $\mathbb{R}$ as common in symplectic homology theory). The cases where the contact Hamiltonian (i.e., the slope) is positive \cite{seidel2008biased} or Reeb-invariant \cite{ritter2016circle} were known prior to \cite{merry2019maximum}. The theory of contact Hamiltonian Floer homology was further developed in \cite{djordjevic2023quantitative, uljarevic2022hamiltonian} and in \cite{cant2023shelukhin, cant2024remarks} with independent contributions.}

\subsection{Immaterial subsets and symplectic homology}

In this section, $W$ denotes a Liouville domain.\footnote{We expect the result of this section to hold in a more general framework than that of Liouville domains. See \cite[Remark~1.5]{uljarevic2023selective}.} Selective symplectic homology $SH_\ast^\Omega(W)$ is a~version of symplectic homology introduced in \cite{uljarevic2023selective}. It is obtained as the direct limit of Floer homology groups for Hamiltonians whose slopes tend to $+\infty$ on the conical end above the open subset $\Omega\subset \partial W $ of the boundary, but remain close to 0 and positive elsewhere. The more precise definition of $SH_\ast^\Omega(W)$ is recalled in Section~\ref{sec:proofSH}.
If $\Omega=\partial W$, then $SH^\Omega_\ast(W)$ coincides with Viterbo's {\cite{viterbo1999functors}} symplectic homology $SH_\ast(W)$. Given two open subsets $\Omega_1\subset \Omega_2$ of $\partial W$, there is a functorial morphism $SH^{\Omega_1}_\ast(W)\to SH^{\Omega_2}_\ast(W)$ called \emph{continuation map}. In particular, the continuation map $SH_\ast^\Omega(W)\to SH_\ast(W)$ is well defined for every open subset $\Omega\subset\partial W$.
The following theorem shows that complements of immaterial Reeb-invariant subsets are big in terms of selective symplectic homology.


\begin{Theorem}\label{thm:complement-big}
Let $A\subset \partial W$ be the preimage of a closed subset under a smooth Reeb-invariant map {$\partial W\to\mathbb{R}$}. Assume $A$ is immaterial and denote $\Omega:=\partial W\setminus A$. Then, the continuation map
\begin{equation*}SH^{\Omega}_\ast(W)\to SH_\ast(W)\end{equation*}
is surjective.
\end{Theorem}


Theorem~\ref{thm:complement-big} can be seen as a generalization and an improvement of \cite[Theorem~7.4]{uljarevic2023selective}. {Namely, \cite[Theorem~7.4]{uljarevic2023selective} claims that the rank of the map $SH^{\Omega}_\ast(W)\to SH_\ast(W)$ is equal to $\dim SH_\ast(W)$ in the case where $A$ is a transversely embedded circle. It is an interesting question whether the condition of $A$ being the preimage of a closed subset under a smooth Reeb-invariant map~${\partial W\to \mathbb{R}}$ can be relaxed.}

\subsection{Contact quasi-measures and immaterial subsets}

In this section, $W$ stands for a weakly\textsuperscript{+} monotone strong filling of a closed contact manifold such that the unit $e\in SH_\ast(W)$ is not eternal. {In what follows, we recall the notions of a weakly\textsuperscript{+} monotone symplectic manifold and of an eternal element. A} $2n$-dimensional symplectic manifold $(W, \omega)$ is weakly\textsuperscript{+} monotone\footnote{ See \cite[Lemma~1.1]{hofer1995floer}, \cite[Assumption ($W^+$)]{seidel1997pi_1}, and \cite[Section~2B]{ritter2016circle}.} if at least one of the following conditions is satisfied:
\begin{enumerate}\itemsep=0pt
\item[(1)] $ \omega |_{\pi_2(W)}=0$,
\item[(2)] $ c_1 |_{\pi_2(W)}=0$,
\item[(3)] there exists $s>0$ such that\footnote{Here, both $\omega$ and $c_1$ are seen as maps assigning a number to an element of $\pi_2(W)$.} $ \omega |_{\pi_2(W)}= s\cdot c_1 |_{\pi_2(W)}$,
\item[(4)] the minimal Chern number of $W$ is at least $n-1$.
\end{enumerate}
These conditions ensure that the symplectic homology of $W$ is well defined. Now, we also recall the notion of being eternal \cite{cant2024remarks}. An element $\theta\in SH_\ast(W)$ is called {\emph{eternal}} if it is contained in the image of the canonical morphism
\[ HF_\ast(H^a)\to SH_\ast(W) \]
for all admissible\footnote{{Recall that a number $a\in\mathbb{R}$ is \emph{admissible} if it is not the period of a closed Reeb orbit on $\partial W$. The extension of this notion to contact Hamiltonians is recalled on page~\pageref{def:admissible}}.} $a\in \mathbb{R}$. Here \smash{$H^a_t\colon \widehat{W}\to \mathbb{R}$} denotes a non-degenerate Hamiltonian on the completion \smash{$\widehat{W}$} of $W$ having slope $a$ at the conical end.

Inspired by work of Entov and Polterovich \cite{Ent14, EP06,EP08,EP-rigidity}, the paper \cite{djordjevic2025quantitative} introduced {a map~$\tau$, called \emph{partial contact quasi-measure},} which assigns a number from $[0,1]$ to a closed subset of~$\partial W$. The quasi-measure $\tau$ is constructed using contact spectral invariants $c(h, \theta)$ introduced independently and simultaneously in \cite{cant2023shelukhin} and \cite{djordjevic2023quantitative}. We recall the construction of $\tau$ in Section~\ref{sec:tau}. The following theorem asserts that each contractible loop of contactomorphisms furnishes a lower bound for~$\tau$.

\begin{Theorem}\label{thm:tau}
Let $W$ and $\tau$ be as above, and $B\subset \partial W$ a closed subset. Let $f\colon\partial W\times [0,1]\to\mathbb{R}$ be a contact Hamiltonian giving rise to a contractible loop of contactomorphisms. Denote {by}~$m$ and $m_B$ the infima of $f$ over $\partial W\times[0,1]$ and $(\partial W\setminus B)\times[0,1]$, respectively. Then, the following inequality holds:
\[\tau(B)\geqslant \frac{m_B}{m_B-m}\]
if $m_B\not=m$.
\end{Theorem}

{By definition, $\tau(B)$ is an element of $[0,1]$. In this paragraph, we explain why Theorem~\ref{thm:tau} does not contradict this fact. The number $m$ cannot be positive, because that would imply that $\partial W$ is not orderable and contradict \cite[Theorem~1.6]{djordjevic2025quantitative}. Therefore, $m_B-m\geqslant m_B$ and $\frac{m_B}{m_B-m}\leqslant 1$. This means that Theorem~\ref{thm:tau} never forces $\tau(B)$ to be greater than 1.

It is useful to observe that Theorem~\ref{thm:tau} is only interesting in the case where $\partial W\setminus B$ is immaterial. Indeed, if $m_B$ is negative, then the theorem asserts that $\tau(B)$ is greater than some negative number, and that follows directly from the definition of $\tau$.}
Lemma~9.3 from \cite{djordjevic2025quantitative} shows $\tau(B)=0$ if $B$ is displaceable and the preimage of a closed subset under a Reeb-invariant function. Therefore, Theorem~\ref{thm:tau} and \cite[Lemma~9.3]{djordjevic2025quantitative} imply the following corollary.


\begin{Corollary}\label{cor:nondisp}
Let $B\subset\partial W$ be the preimage of a closed subset under a Reeb-invariant smooth function {$\partial W\to\mathbb{R}$}. If $\overline{\partial W \setminus B}$ is immaterial, then $B$ is contact non-displaceable.
\end{Corollary}

It is an interesting question whether the condition of $B$ being the preimage of a closed subset under a smooth Reeb-invariant map $\partial W\to \mathbb{R}$ can be relaxed.
The next example applies Theorem~\ref{thm:tau} in a concrete situation of $S^\ast \mathbb{S}^n$. The example illustrates a general method of finding immaterial subsets in unit cotangent bundles by lifting contractible loops of diffeomorphisms from the base.

\begin{Example}
Let $n\in\mathbb{N}$ be a number equal to 3 modulo 4 and let $\psi_t\colon \mathbb{S}^n\to\mathbb{S}^n$ be given by $\psi_t(x):= {\rm e}^{2\pi {\rm i} t}\cdot x$ (we see $\mathbb{S}^n$ as the unit sphere in the complex space). Since $n\geqslant 3$ is equal to 3 modulo 4, $\psi_t$ is a contractible loop of diffeomorphisms.\footnote{The loop $\psi_t$ represents an even power of the generator of $\pi_1\bigl(U\bigl(\frac{n+1}{2}\bigr)\bigr)\cong\mathbb{Z}$ because $\frac{n+1}{2}$ is even. Therefore, $\psi_t$ represents the trivial element in $\pi_1(O(n+1))\cong\mathbb{Z}_2$. Consequently, $\psi_t$ is contractible as a loop of diffeomorphisms.} Therefore,
\[
\varphi_t\colon \ S^\ast\mathbb{S}^n\to S^\ast\mathbb{S}^n, \qquad v^\ast\mapsto \psi^\ast_t(v^\ast)
\]
is a contractible loop of contactomorphisms. The contact Hamiltonian of $\varphi_t$ is given by
\[
f\colon \ S^\ast\mathbb{S}^n\to[-2\pi, 2\pi], \qquad f(v^\ast):= v^\ast(X(p)),
\]
where $X$ is the vector field of $\psi_t$ and $p$ is the base point for $v^\ast$. Denote by $B_c$ the subset $\{v^\ast\in S^\ast\mathbb{S}^n\mid v^\ast(X)\leqslant c\}$. The set $B_c$ for $c\in(-2\pi, 2\pi), $ can be seen as a fibration over $\mathbb{S}^n$ with each fibre diffeomorphic to the $(n-1)$-dimensional ball $\mathbb{B}^{n-1}$. Now, Theorem~\ref{thm:tau} implies
\[ \tau(B_c)\geqslant \frac{c}{c+2\pi}. \]
As a consequence of Corollary~\ref{cor:nondisp}, the set $B_c$ is contactly non-displaceable for $c>0$. This claim, however, is not new. In fact, \cite[Example~2.11]{sun2026contact} asserts that $B_c$ is non-displaceable for $c\geqslant 0$.
\end{Example}


\section{Proof of Theorem~\ref{thm:complement-big}}\label{sec:proofSH}

 We first recall the definition of the selective symplectic homology $SH_\ast^\Omega(W)$. For an open subset $\Omega\subset \partial W$, denote by $\mathcal{H}_\Omega$ the set of the smooth autonomous contact Hamiltonians $h\colon \partial W\to[0,+\infty)$ such that\footnote{The second condition is slightly different from that of \cite[Definition~3.1]{uljarevic2023selective}. However, using either of these two conditions yields the same definition of selective symplectic homology.}
 \begin{enumerate}\itemsep=0pt
 \item[(1)] ${\rm supp} h \subset \Omega$,
 \item[(2)] $h$ vanishes up to order $2$ at each $p\in\partial W$ such that $h(p)=0$,
 \item[(3)] the 1-periodic orbits of $h$ are constant.
 \end{enumerate}

 Denote by $\Pi(h)$ the set of the positive smooth functions $f\colon \partial W\to\mathbb{R}_{{>}0}$ such that the contact Hamiltonian $h+f$ has no 1-periodic orbits. {We see $\mathcal{H}_\Omega$ and $\Pi(h)$ as directed sets with the preorder relation pointwise $\leqslant$.} The selective symplectic homology is defined by
 \begin{equation}\label{eq:SSH} SH_\ast^\Omega(W):= \underset{h\in\mathcal{H}_\Omega}{\lim_{\longrightarrow}} \underset{f\in\Pi(h)}{\lim_{\longleftarrow}} HF_\ast(h+f),\end{equation}
 where the limits are taken with respect to the continuation maps {\cite[Section~4]{merry2019maximum}}. Here, $HF_\ast(h)$ stands for the Floer homology for the contact Hamiltonian $h$, that was introduced in \cite{merry2019maximum}. Notice that the inverse limit in \eqref{eq:SSH} stabilizes \cite[Lemma~3.5]{uljarevic2023selective}. In the following lemma and in the rest of the paper, \smash{$\varphi^f_t$} denotes the contact isotopy furnished by the contact Hamiltonian $f\colon\partial W\times\mathbb{R}\to\mathbb{R}$. {The contact Hamiltonian of the composition \smash{$\varphi_t^h\circ\varphi_t^f$} is denoted by $h\# f$. Explicitly,
 \[ (h\# f)_t(x):= h_t(x) + \bigl(\kappa_t^h\cdot f_t\bigr)\circ\bigl(\varphi^h_t\bigr)^{-1}(x), \]
 where $\kappa^h_t$ is defined by \smash{$\bigl(\varphi^h_t\bigr)^\ast\alpha =\kappa_t^h\cdot \alpha$}.}

\begin{Lemma}\label{lem:twisttrick}
In the situation of Theorem~$\ref{thm:complement-big}$, for every $\ell\in\mathbb{R}$ there exists $h\in\mathcal{H}_{\Omega}(\partial W)$ and a~contractible loop \smash{$\varphi^f_t\colon \partial W\to\partial W$} of contactomorphisms\footnote{Here and in the rest of the paper, loops $\varphi_t\colon \partial W\to\partial W$ of contactomorphisms are seen as 1-periodic smooth $\mathbb{R}$-families $\{\varphi_t\}_{t\in \mathbb{R}}$ of contactomorphisms.} such that
\[ (h+\varepsilon)\# f\geqslant \ell\]
for all $\varepsilon \geqslant 0$.
\end{Lemma}
\begin{proof}
Without loss of generality, assume $\ell>0$. By the assumptions in Theorem~\ref{thm:complement-big}, {the set $A$ is the preimage of a closed subset under a smooth Reeb-invariant map $\partial W\to\mathbb{R}$. Since any closed subset of $\mathbb{R}$ can be seen as the preimage of $\{0\}$ under some smooth function $\mathbb{R}\to\mathbb{R}$,} there exists a smooth Reeb-invariant function\footnote{{Obtained as the composition.}} $\mu\colon \partial W\to\mathbb{R}$ such that $A=\mu^{-1}(\{0\})$. Since $A$ is compact and immaterial, there exists (see \cite[Lemma~7.2]{uljarevic2023selective}) a contractible loop of contactomorphisms \smash{$\varphi^f_t\colon \partial W\to\partial W$} such that
\begin{equation} \label{eq:main2ell}
\min_{x\in A, t\in\mathbb{R}} f_t(x)\geqslant 2\ell.
\end{equation}
The set
\[S:=\bigl\{ x\in\partial W\mid \bigl(\min_{t\in\mathbb{R}} f_t\bigr)(x)\leqslant \ell \bigr\}\]
is closed because the function
\[ \min_{t\in\mathbb{R}} f_t = \min_{t\in[0,1]} f_t \colon \ \partial W\to \mathbb{R} \]
is {continuous}. In addition, $S$ and $A$ are disjoint due to \eqref{eq:main2ell}. Hence, there exists $a>0$ such that $\mu^{-1}(-a, a)$ is disjoint from $S$. Let $m:=\min f$ be the (global) minimum of $f$ and let $\beta\colon \mathbb{R}\to[0,+\infty)$ be a smooth function such that
\begin{enumerate}\itemsep=0pt
\item[(1)] there exists $\delta>0$ such that $\beta(s)=0$ if, and only if, $s\in[-\delta, \delta]$, and
\item[(2)] $\beta(s)\geqslant \ell - m$ for $s\not\in (-a, a)$.
\end{enumerate}
Consequently, $\beta\circ \mu(x)\geqslant \ell-m$ for $x$ such that $\min_{t\in\mathbb{R}} f_t(x)\leqslant \ell$ and ${\rm supp} \beta\circ \mu \subset \Omega= \partial W\setminus A$. Denote $h:=\beta\circ\mu$. By the construction of $h$, we have $h+f_t\geqslant \ell$. Therefore,
\begin{align*}
((h + \varepsilon)\# f)_t(x) &= h(x) +\varepsilon + f_t\circ\bigl(\varphi^{h+\varepsilon}_t\bigr)^{-1}(x) = (h +\varepsilon + f_t)\circ\bigl(\varphi^{h+\varepsilon}_t\bigr)^{-1}(x) \geqslant \ell.
\end{align*}
Here, we used \smash{$ (h +\varepsilon )\circ\bigl(\varphi^{h+\varepsilon}_t\bigr)^{-1}= h + \varepsilon$}, which holds because $h+\varepsilon$ is time-independent and Reeb-invariant. By perturbing $h$ slightly {(within the class of non-negative smooth functions having the same support as $h$)} if necessary, one can achieve $h\in\mathcal{H}_\Omega$. This completes the proof.
\end{proof}

A contact Hamiltonian $h$ is called \emph{admissible}\label{def:admissible} if its time-1 map $\varphi^h_1$ has no fixed points $x$ such that \smash{$\bigl(\bigl(\varphi^h_1\bigr)^\ast\alpha\bigr)\bigr|_x= \alpha |_x$}. The admissibility condition is required for $HF_\ast(h)$ to be well defined. In the special case of a constant contact Hamiltonian $h\equiv a\in \mathbb{R}$, the admissibility simply means that $a$ is not a period of a Reeb orbit.
The following lemma asserts that the images of $HF_\ast(h)$ and $HF_\ast(g)$ under the canonical morphisms coincide provided $h$ and $g$ give rise to the same element in \smash{$\widetilde{{\rm Cont}}_0(\partial W)$}.

\begin{Lemma}\label{lem:compwithcontractible}
Let $h_t\colon \partial W\to \mathbb{R}$ be an admissible contact Hamiltonian and let $\varphi^f_t\colon \partial W\to\partial W$ be a contractible loop of contactomorphisms. Then, the images of the following natural morphisms are equal:
\begin{align*}
& HF_\ast(h)\to SH_\ast(W),\\
& HF_\ast(h\#f)\to SH_\ast(W).
\end{align*}
\end{Lemma}
\begin{proof}
Since $f$ gives rise to a contractible loop of contactomorphisms, there exists a {\emph{zig-zag isomorphism}} $HF_\ast(h)\to HF_\ast(h\# f)$ (see \cite[Section~9]{uljarevic2023selective} or \cite[Section~2.1.3]{djordjevic2025quantitative}). Proposition~2.4 from \cite{djordjevic2025quantitative} asserts that the following diagram, consisting of the canonical morphisms and a zig-zag isomorphism, commutes:
\[\begin{tikzcd}
&SH_\ast(W)&\\
HF_\ast(h)\arrow["\text{zig-zag}"]{rr}\arrow{ru}&&HF_\ast(h\#f).\arrow{lu}
\end{tikzcd}\]
From this commutative diagram, the claim of the lemma follows directly.
\end{proof}

\begin{proof}[Proof of Theorem~\ref{thm:complement-big}]
For an admissible contact Hamiltonian $h_t\colon \partial W\to\mathbb{R}$, denote by $\mathfrak{Im}(h)$ the image of the natural morphism $HF_\ast(h)\to SH_\ast(W)$. Let $\theta\in SH_\ast(W)$. Then, there exists a~real number $\ell\in \mathbb{R}$ that is not a period of a Reeb orbit on $\partial W$ such that $\theta\in\mathfrak{Im}(\ell)$. By Lemma~\ref{lem:twisttrick}, there exist $h\in\mathcal{H}_\Omega(\partial W)$ and a contractible loop $\varphi^f_t\colon \partial W\to\partial W$ of contactomorphisms such that the natural morphism $HF_\ast(\ell)\to SH_\ast(W)$ factors through $HF_\ast((h+\varepsilon)\#f)\to SH_\ast(W)$ for all positive $\varepsilon \in \Pi(h)$. Hence, $\mathfrak{Im}(\ell)\subset \mathfrak{Im}((h+\varepsilon)\# f)$ for all positive $\varepsilon \in \Pi(h)$. Since $f$ generates a contractible loop of contactomorphisms, Lemma~\ref{lem:compwithcontractible} implies $\mathfrak{Im}(\ell)\subset \mathfrak{Im}(h+\varepsilon)$ for all positive $\varepsilon \in \Pi(h)$. This further implies that $\theta$ belongs to the image of the continuation map $SH_\ast^\Omega(W)\to SH_\ast(W)$. Since $\theta$ was arbitrarily chosen, this completes the proof.
\end{proof}



\section{Proof of Theorem~\ref{thm:tau}}\label{sec:tau}

The quasi-measure $\tau$ is constructed in \cite{djordjevic2025quantitative} using contact spectral invariants $c(h, \theta)$ introduced independently and simultaneously in \cite{cant2023shelukhin} and \cite{djordjevic2023quantitative}.\footnote{In the present paper, we follow the approach and the conventions from \cite{djordjevic2025quantitative}} The contact spectral invariant $c(h, \theta)$ is assigned to an element $\theta\in SH_\ast(W)$ and a time-dependent contact Hamiltonian $h_t\colon \partial W\to \mathbb{R}$. {For} the construction of $\tau$, the case where $\theta$ is equal to the unit $e\in SH_\ast(W)$ is relevant.
Now, we recall the properties of the contact spectral invariants that are needed in the present paper:
\begin{enumerate}\itemsep=0pt
\item[(1)] $c(s, e)= s$ for all $s\in \mathbb{R}$.
\item[(2)] ({Invariance}) If $\bigl\{\varphi^h_t\bigr\}$ and $ \{\varphi_t^g \}$ represent the same class in \smash{$\widetilde{{\rm Cont}}_0(\partial W)$}, then $c(h,e)= c(g, e)$.
\item[(3)] (Monotonicity) If $h\leqslant g$ pointwise, then ${c(h, e)\leqslant c(g, e)}$.
\item[(4)] (Triangle inequality) If $h$ and $g$ are Reeb invariant, then\footnote{Notice that, due to our conventions, the triangle inequality is ``reversed''.}
\[
c(h, e) + c(g, e)\leqslant c(h\#g, e).
\]
\item[(5)] (Stability) If $h$ and $g$ are Reeb-invariant, then
\[ \min (h-g)\leqslant c(h,e) - c(g, e)\leqslant \max(h-g). \]
\end{enumerate}

For a Reeb-invariant contact Hamiltonian $h\colon \partial W\to\mathbb{R}$, the contact quasi-state $\zeta(h)$ is defined~by
\[\zeta(h):= \lim_{k\to+\infty} \frac{c(k\cdot h, e)}{k}.\]
The limit exists and is finite \cite[Lemma~7.2]{djordjevic2025quantitative}. The contact quasi-state $\zeta$ satisfies the following homogeneity property: $\zeta(sh)= s \zeta(h)$ for $s\geqslant 0$. Now, the contact quasi-measure $\tau$ is defined by
\[ \tau(A):= \inf \{ \zeta(h)\mid h\colon\partial W\to[0,1] \text{ Reeb-invariant, smooth, and } h|_A=1 \} \]
for a closed subset $A\subset\partial W$. The following lemma is used in the proof of Theorem~\ref{thm:tau}.

{
\begin{Lemma}\label{lem:aut}
Let $M$ be a closed contact manifold with a~contact form $\alpha$ and let $B\subset M$ be a~closed subset. Let $f_t\colon M\to\mathbb{R}$ be a~contact Hamiltonian and $h\colon M\to[0,1]$ a Reeb-invariant contact Hamiltonian such that $ h |_B=1$. Denote
\[\ell:=\inf_{t\in \mathbb{R}, x\in M\setminus B} f_t(x)\qquad \text{and} \qquad
m:=\inf_{t\in \mathbb{R}, x\in M} f_t(x),\]
and assume $m$ is finite.\footnote{This implies $\ell$ is finite as well.} Then, the following inequality holds:
\[ ((\ell-m)\cdot h)\# f\geqslant \ell. \]
\end{Lemma}}
\begin{proof}
Denote $g:=(\ell-m)\cdot h$. Since $g$ is time-independent and Reeb-invariant, it is constant along its trajectories. Consequently,
\begin{align*}
(g\# f)_t(x) = g(x) + f_t\circ (\varphi^{g}_t )^{-1}(x) = (g + f_t)\circ (\varphi^{g}_t )^{-1}(x).
\end{align*}
Therefore, it is enough to prove $g+f_t\geqslant \ell$. If $x\in {M\setminus B}$, then $f_t(x)\geqslant \ell$. Since $g$ is non-negative, this implies $g(x)+ f_t(x)\geqslant \ell$. If, on the other hand, $x\not \in {M\setminus B}$, then $x\in B$ and
\[ g(x) + f_t(x)= \ell-m + f_t(x)\geqslant \ell. \]
This finishes the proof.
\end{proof}

\begin{proof}[Proof of Theorem~\ref{thm:tau}]
Let $h\colon\partial W\to[0,1]$ be a Reeb-invariant contact Hamiltonian such that~${ h |_B=1}$. By Lemma~\ref{lem:aut}, we get $ ((m_B-m)\cdot h )\#f\geqslant m_B$. The {invariance} and the monotonicity properties of the contact spectral invariants imply
\[c((m_B-m)\cdot h, e)= c ( ((m_B-m)\cdot h )\#f, e )\geqslant c(m_B, e)=m_B.\]
The same argument applied to the iterated loop $\varphi^f_{kt}$ (generated by the contact Hamiltonian ${k\cdot f_{kt}}$), $k\in\mathbb{N}$, implies
\[c(k\cdot(m_B-m)\cdot h, e)\geqslant k\cdot m_B.\]
Therefore,
\[\zeta((m_B-m)\cdot h)= \lim_{k\to+\infty} \frac{c(k\cdot(m_B-m)\cdot h, e)}{k}\geqslant m_B. \]
Consequently, due to homogeneity of $\zeta$, we have \smash{$\zeta(h)\geqslant\frac{m_B}{m_B-m}$}. This completes the proof.
\end{proof}

\begin{thebibliography}{99}
\footnotesize\itemsep=0pt

\bibitem{cant2023shelukhin}
Cant D., Shelukhin's Hofer distance and a~symplectic cohomology barcode for
 contactomorphisms,
 \href{https://doi.org/10.1142/S179352532650010X}{\textit{J.~Topol. Anal.}}, {t}o appear, \href{http://arxiv.org/abs/2309.00529}{arXiv:2309.00529}.

\bibitem{cant2024remarks}
Cant D., Remarks on eternal classes in symplectic cohomology,
 \href{http://arxiv.org/abs/2410.03914}{arXiv:2410.03914}.


\bibitem{djordjevic2023quantitative}
Djordjevi\'c D., Uljarevi\'c I., Zhang J., Quantitative characterization in
 contact {H}amiltonian dynamics -- {I},
 \href{http://arxiv.org/abs/2309.00527}{arXiv:2309.00527}.

\bibitem{djordjevic2025quantitative}
Djordjevi\'c D., Uljarevi\'c I., Zhang J., Quantitative contact {H}amiltonian
 dynamics, \href{http://arxiv.org/abs/2507.13234}{arXiv:2507.13234}.

\bibitem{eliashberg2006geometry}
Eliashberg Y., Kim S.S., Polterovich L., Geometry of contact transformations
 and domains: orderability versus squeezing,
 \href{https://doi.org/10.2140/gt.2006.10.1635}{\textit{Geom. Topol.}}
 \textbf{10} (2006), 1635--1747,
 \href{http://arxiv.org/abs/math.GR/0511658}{arXiv:math.GR/0511658}.

\bibitem{eliashberg2000partially}
Eliashberg Y., Polterovich L., Partially ordered groups and geometry of contact
 transformations, \href{https://doi.org/10.1007/PL00001656}{\textit{Geom.
 Funct. Anal.}} \textbf{10} (2000), 1448--1476,
 \href{http://arxiv.org/abs/math.SG/9910065}{arXiv:math.SG/9910065}.

\bibitem{Ent14}
Entov M., Quasi-morphisms and quasi-states in symplectic topology, in
 Proceedings of the {I}nternational {C}ongress of {M}athematicians~-- {S}eoul
 2014. {V}ol.~{II}, Kyung Moon Sa, Seoul, 2014, 1147--1171,
 \href{http://arxiv.org/abs/1404.6408}{arXiv:1404.6408}.

\bibitem{EP06}
Entov M., Polterovich L., Quasi-states and symplectic intersections,
 \href{https://doi.org/10.4171/CMH/43}{\textit{Comment. Math. Helv.}}
 \textbf{81} (2006), 75--99,
 \href{http://arxiv.org/abs/math.SG/0410338}{arXiv:math.SG/0410338}.

\bibitem{EP08}
Entov M., Polterovich L., Symplectic quasi-states and semi-simplicity of
 quantum homology, in Toric Topology, \textit{Contemp. Math.}, Vol.~460,
American Mathematical Society,
 Providence, RI, 2008, 47--70,
 \href{http://arxiv.org/abs/0705.3735}{arXiv:0705.3735}.

\bibitem{EP-rigidity}
Entov M., Polterovich L., Rigid subsets of symplectic manifolds,
\textit{Compos. Math.}
 \textbf{145} (2009), 773--826,
 \href{http://arxiv.org/abs/0704.0105}{arXiv:0704.0105}.

\bibitem{hofer1995floer}
Hofer H., Salamon D.A., Floer homology and {N}ovikov rings, in The {F}loer
 Memorial Volume, \textit{Progr. Math.}, Vol.~133,
 \href{https://doi.org/10.1007/978-3-0348-9217-9_20}{Birkh\"auser}, Basel,
 1995, 483--524.

\bibitem{merry2019maximum}
Merry W.J., Uljarevic I., Maximum principles in symplectic homology,
 \href{https://doi.org/10.1007/s11856-018-1792-z}{\textit{Israel~J. Math.}}
 \textbf{229} (2019), 39--65,
 \href{http://arxiv.org/abs/1705.06108}{arXiv:1705.06108}.

\bibitem{oh2004construction}
Oh Y.-G., Construction of spectral invariants of {H}amiltonian paths on closed
 symplectic manifolds, in The Breadth of Symplectic and {P}oisson Geometry,
 \textit{Progr. Math.}, Vol.~232,
 \href{https://doi.org/10.1007/0-8176-4419-9_18}{Birkh\"auser}, Boston, MA,
 2005, 525--570,
 \href{http://arxiv.org/abs/math.SG/0405064}{arXiv:math.SG/0405064}.

\bibitem{polterovich2014function}
Polterovich L., Rosen D., Function theory on symplectic manifolds,
 \textit{{CRM} Monogr. Ser.}, Vol.~34,
 \href{https://doi.org/10.1090/crmm/034}{American Mathematical Society},
 Providence, RI, 2014.

\bibitem{ritter2016circle}
Ritter A.F., Circle actions, quantum cohomology, and the {F}ukaya category of
 {F}ano toric varieties,
 \href{https://doi.org/10.2140/gt.2016.20.1941}{\textit{Geom. Topol.}}
 \textbf{20} (2016), 1941--2052,
 \href{http://arxiv.org/abs/1406.7174}{arXiv:1406.7174}.

\bibitem{seidel1997pi_1}
Seidel P., {$\pi_1$} of symplectic automorphism groups and invertibles in
 quantum homology rings,
 \href{https://doi.org/10.1007/s000390050037}{\textit{Geom. Funct. Anal.}}
 \textbf{7} (1997), 1046--1095,
 \href{http://arxiv.org/abs/dg-ga/9511011}{arXiv:dg-ga/9511011}.

\bibitem{seidel2008biased}
Seidel P., A~biased view of symplectic cohomology, in Current Developments in
 Mathematics,
 \href{https://doi.org/10.4310/CDM.2006.v2006.n1.a4}{International Press},
 Somerville, MA, 2008, 211--253,
 \href{http://arxiv.org/abs/0704.2055}{arXiv:0704.2055}.

\bibitem{sun2026contact}
Sun Y., Uljarevi\'c I., Varolgunes U., Contact big fiber theorems,
 \href{https://doi.org/10.1007/s00039-026-00734-4}{\textit{Geom. Funct.
 Anal.}} \textbf{36} (2026), 301--350,
 \href{http://arxiv.org/abs/2503.04277}{arXiv:2503.04277}.

\bibitem{uljarevic2023selective}
Uljarevi\'c I., Selective symplectic homology with applications to contact
 non-squeezing,
 \href{https://doi.org/10.1112/s0010437x23007480}{\textit{Compos. Math.}}
 \textbf{159} (2023), 2458--2482,
 \href{http://arxiv.org/abs/2205.14771}{arXiv:2205.14771}.

\bibitem{uljarevic2022hamiltonian}
Uljarevi\'c I., Zhang J., Hamiltonian perturbations in contact {F}loer
 homology, \href{https://doi.org/10.1007/s11784-022-00986-1}{\textit{J.~Fixed
 Point Theory Appl.}} \textbf{24} (2022), 71, 20~pages,
 \href{http://arxiv.org/abs/2203.12500}{arXiv:2203.12500}.

\bibitem{viterbo1999functors}
Viterbo C., Functors and computations in {F}loer homology with
 applications.~{I}, \href{https://doi.org/10.1007/s000390050106}{\textit{Geom.
 Funct. Anal.}} \textbf{9} (1999), 985--1033.

\end{thebibliography}
\end{document}
